\documentclass[12pt]{article}
\usepackage{notes-project}

\begin{document}

\begin{center}
  {\LARGE\bfseries Lecture 2: Cardinality}\\[0.25em]
  {\large Mathematical Analysis}\\[0.15em]
\end{center}

\section{Nested Sequence of Sets} 

\begin{dfn}[Nested Sequence of Sets]{nested-sequence-of-sets}
  A sequence of sets \((A_n)_{n\in\NN}\) is \emph{nested} if
  \[
    \forall n\in\NN,\quad A_{n+1}\subseteq A_n.
  \] 
\end{dfn}


\begin{thm}[Nested Sequence Containment]{nested-sequence-containment}
  Let \((A_n)_{n\in\NN}\) be a
  \xref[nested sequence of sets]{def:nested-sequence-of-sets}. Then
  \[
    \forall p,q\in\NN,\qquad
    p\leq q\implies A_q\subseteq A_p.
  \]
\end{thm}

\begin{proof}
\color{blue}\textbf{TODO} \color{black}
\end{proof}

\section{Intervals}

\begin{dfn}[Interval]{interval}
  A set \(I\subseteq\RR\) is an \emph{interval} if, whenever
  \(x,y\in I\) with \(x\le y\), every \(z\in\RR\) satisfying
  \(x\le z\le y\) also belongs to \(I\).
\end{dfn}

\begin{rem}[Intervals as Sets]{intervals-as-sets}
  In these notes, every interval is a set contained in \(\RR\);
  that is, every interval \(I\) satisfies \(I\subseteq\RR\).
\end{rem}

\begin{dfn}[Open Interval]{open-interval}
  For \(a,b\in\RR\) with \(a<b\), the \emph{open interval}
  from \(a\) to \(b\) is
  \[
    (a,b)=\{x\in\RR:a<x<b\}.
  \]
\end{dfn}

\begin{dfn}[Closed Interval]{closed-interval}
  For \(a,b\in\RR\) with \(a\le b\), the \emph{closed interval}
  from \(a\) to \(b\) is
  \[
    [a,b]=\{x\in\RR:a\le x\le b\}.
  \]
\end{dfn}

\begin{lem}[Endpoints of Contained Closed Intervals]{contained-interval-endpoints}
  Let \(I_1=[a_1,b_1]\) and \(I_2=[a_2,b_2]\) be closed intervals. If
  \(I_1\subseteq I_2\), then
  \[
    a_2\le a_1\le b_1\le b_2.
  \]
\end{lem}

\begin{proof}
  Since \(a_1,b_1\in I_1\subseteq I_2\),
  \xref{def:closed-interval} gives \(a_2\le a_1\) and \(b_1\le b_2\).

  Since \(I_1=[a_1,b_1]\) is a closed interval,
  \xref{def:closed-interval} requires \(a_1\le b_1\).

  Combining these inequalities gives the result.
\end{proof}

\begin{lem}[Endpoint Order of Nested Closed Intervals]{nested-interval-endpoint-order}
  Let \(I_k=[a_k,b_k]\) be a nested sequence of closed intervals. Then
  \[
    n\leq m
    \implies
    a_n\leq a_m\leq b_m\leq b_n.
  \]
\end{lem}

\begin{proof}
\color{blue}\textbf{TODO} \color{black}
\end{proof}

\begin{thm}[Nested Intervals Property (NIP)]{nip}
  Let \(I_n=[a_n,b_n]\) be a sequence of
  \xref[closed]{def:closed-interval} and \xref[bounded]{lec01::def:bounded-set} \xref[intervals]{def:interval}.
  \[
    I_0\supseteq I_1\supseteq I_2\supseteq\cdots
    \implies
    \bigcap_{n=0}^{\infty} I_n\neq\emptyset.
  \]
\end{thm}

\begin{proof}
Let \(A=\{a_n:n\in\NN\}\). The nesting hypothesis gives
\[
  \forall n\in\NN,\quad I_{n+1}\subseteq I_n.
\]

For fixed \(p\in\NN\), define
\[
  P_p(r)\iff I_{p+r}\subseteq I_p.
\]
\[
  P_p(0),\quad
  P_p(r)\implies
  I_{p+r+1}\subseteq I_{p+r}\subseteq I_p
  \implies P_p(r+1).
\]
By \xref[Natural Number Induction]{lec01::thm:natural-number-induction}
and \(p\leq q\iff\exists r\in\NN:q=p+r\),
\[
  \forall p,q\in\NN,\quad
  p\leq q\implies
  \exists r\in\NN:\quad
  q=p+r\ \land\ I_q=I_{p+r}\subseteq I_p.
\]

By \xref[comparability in an ordered set]{lec01::def:ordered-set},
\(\forall n,m\in\NN,\ (n\leq m)\lor(m\leq n)\). Hence
\xref{lem:contained-interval-endpoints} gives
\[
  \begin{aligned}
    n\leq m &\implies I_m\subseteq I_n
      \implies a_n\leq a_m\leq b_m,\\
    m\leq n &\implies I_n\subseteq I_m
      \implies a_n\leq b_n\leq b_m,
  \end{aligned}
\]
so
\[
  \forall n,m\in\NN,\quad a_n\leq b_m.
\] 

Consequently,
\[
  a_0\in A,\, 
  \forall n\in\NN,\ a_n\leq b_0
  \implies
  A\neq\emptyset\ \land A\text{ is bounded above}.
\]
By the \xref[Axiom of Completeness]{lec01::axm:aoc}, choose
\(\alpha=\sup A\). The
\xref[definition of supremum]{lec01::def:supremum} gives
\[
  \forall n,m\in\NN,\quad
  a_n\leq\alpha\leq b_m.
\]
Hence, by the
\xref[definitions of a closed interval]{def:closed-interval} and
\xref[indexed intersection]{lec01::def:indexed-intersection},
\[
  \left(\forall n\in\NN,\quad a_n\leq\alpha\leq b_n\right)
  \implies
  \left(\forall n\in\NN,\quad \alpha\in I_n\right)
  \iff
  \alpha\in\bigcap_{n=0}^{\infty}I_n.
\]
Therefore \(\bigcap_{n=0}^{\infty}I_n\neq\emptyset\).
\end{proof}

\begin{rem}[Intervals can be Infinitely Small]{inf-small-intervals}
If a size of the interval is infinitely small, ie \(I=[a,b]\),
\(\forall\varepsilon>0,\,
b-a<\varepsilon\), this means that the interval is equivalent to a single point, ie $a = b$. 
\end{rem}

\section{Decimal Expansion} 

\begin{thm}[Decimal Expansion]{decimal-representation}
Let \(D=\{0,1,\ldots,9\}\). Every digit sequence
\(a:\NNOne\to D\) determines a unique \(x\in[0,1]\):
\[
  \forall a:\NNOne\to D,\quad
  \exists!x\in[0,1],\qquad
  x=\sum_{k=1}^{\infty}a_k10^{-k}.
\]
\end{thm}

\begin{proof}
We call the summation of the first $n$ decimals $s_n$, ie \[s_n  = \sum^{n}_{k=1} a_k \, 10^{-k}\]

Since \(\forall k \in \NNOne, 0 \leq a_k \leq 9\), for all
\(n,m\in\NNOne\) with \(n\leq m\),
\begin{align*}
  s_m
  &= s_n + \sum_{k=n+1}^{m} a_k 10^{-k} \\
  &\leq s_n + 9 \sum_{k=n+1}^{m} 10^{-k} \\
  &= s_n + 10^{-n} - 10^{-m}
  < s_n + 10^{-n}.
\end{align*}

In other words, \[\forall n,m \in \NNOne,\quad n \leq m \implies s_m \in [s_n, s_n + 10^{-n}]\] 

By the same argument, we can also show that \[\forall n,m \in \NNOne,\quad n \leq m \implies s_m + 10^{-m} \in [s_n, s_n + 10^{-n}]\] 

Where we can construct the following nested intervals
\[
I_n=[s_n,s_n+10^{-n}],\qquad
I_1\supseteq I_2\supseteq\cdots, \qquad x = \sum_{n = \NNOne }a_n 10^{-k}= \bigcap_{n\in\NNOne}I_n
\]

According to the \xref[Nested Interval Theorem]{thm:nip}, after reindexing,
\[
\bigcap_{n\in\NNOne}I_n\neq\emptyset.
\]

Also, notice that $I_n = [s_n, s_n + 10^{-n}]$, $s_n + 10^{-n} - s_n = 10^{-n}$. According to the \xref[Remark]{rem:inf-small-intervals}, we can conclude:
\[
\forall\varepsilon>0,\ \exists n\in\NNOne:\,
I_n=[s_n,s_n+10^{-n}]
\,\text{and}\,
(s_n+10^{-n})-s_n=10^{-n}<\varepsilon \implies \exists!x\in\bigcap_{n\in\NNOne}I_n
\]

Thus \(x\) is the unique number represented by the digit sequence \(a\).
\end{proof} 

\begin{rem}[Decimal Expansion is not Unique]{decimal-representation-not-unique}
  The decimal expansion of a real number is not unique. For example, \(0.8999\ldots = 0.9000\ldots\)  
\end{rem}

\begin{thm}[Uniqueness of Canonical Decimal Expansions]{canonical-decimal-uniqueness}
Let \(D=\{0,1,\ldots,9\}\). For every \(x\in[0,1)\), there exists a unique
digit sequence \(a:\NNOne\to D\) which is not eventually equal to \(9\) and
satisfies
\[
  x=\sum_{n=1}^{\infty}a(n)10^{-n}.
\]
Here, ``not eventually equal to \(9\)'' means
\[
  \neg\exists N\in\NNOne\;\forall n\geq N,\quad a(n)=9.
\]
\end{thm}

\begin{proof}
\color{blue}\textbf{TODO} \color{black}
\end{proof}

\section{Archimedean Property} 

\begin{thm}[Archimedean Property]{archimedean-property}
  $\NN$ is not \xref[bounded above]{lec01::def:bounded-above} in $\RR$, in other words, \(\forall x \in \RR, \exists n \in \NN, n > x\)
\end{thm} 

\begin{proof}
Assume for the sake of contradiction that $\NN$ is bounded above in $\RR$. 

Since \(0\in\NN\), we have \(\NN\neq\emptyset\). By the \xref[Axiom of Completeness]{lec01::axm:aoc}, $\NN$ has a supremum, denoted as $\alpha = \sup \NN$.

By the definition of \xref[epsilon criteria]{lec01::prop:epsilon-supremum}, we have: 
\[\forall x \in \NN, x \leq \alpha, \quad \forall \varepsilon > 0, \exists x \in \NN, \alpha - \varepsilon < x \] 

Taking \(\varepsilon=1\), we obtain
\[\exists n\in\NN:\quad \alpha-1<n\implies\alpha<n+1.\]

Since \(n+1=S(n)\in\NN\), this contradicts that \(\alpha\) is an upper bound of \(\NN\), namely
\[\forall m\in\NN,\quad m\leq\alpha.\]
Therefore, $\NN$ is not bounded above in $\RR$. 
\end{proof}

\begin{thm}[Archimedean Property (Reciprocal Form)]{archimedean-reciprocal-form}
\[\forall y > 0, \, \exists n \in \NN, \, 1/n < y \]
\end{thm}

\begin{proof}
Take $x = 1/y$. Since $y > 0$, we have $x > 0$. 

According to definition of Natural Numbers, $\exists n \in \NN, \, n \geq x$, then $1/n \leq 1/x = y$. 
\end{proof}

\begin{dfn}[Dense Subset]{dense-subset}
Let \(D\subseteq\RR\). We say that \(D\) is \emph{dense in \(\RR\)} if
\[
\forall a,b\in\RR,\quad
a<b\implies\exists d\in D:\quad a<d<b.
\]
\end{dfn}

\begin{thm}[Density of the Rational Numbers]{rationals-dense-in-reals}
  $\QQ$ is dense in $\RR$. 
\end{thm}

\begin{proof}
\color{blue}\textbf{TODO} \color{black} 
\end{proof}

\begin{thm}[Density of the Irrational Numbers]{irrationals-dense-in-reals}
  $\RR \setminus \QQ$ (ie, $\II$) is dense in $\RR$. 
\end{thm} 

\begin{proof}
\color{blue}\textbf{TODO} \color{black} 
\end{proof} 

\begin{rem}
$\QQ$ and $\II$ are both dense in $\RR$. It does not mean that $\QQ$ and $\II$ are having the same size. The size of infinite sets are defined in 
\end{rem}

\section{Cardinality} 

\begin{dfn}[Cardinality]{cardinality}
  The \emph{cardinality} of a set \(A\), denoted by \(\card{A}\), is the number of elements in \(A\).
\end{dfn}

\begin{dfn}[Equinumerous Sets]{same-cardinality}
Sets \(A\) and \(B\) are \emph{equinumerous}, written \(A\equinum B\), if and only if there exists a
\xref[bijective mapping]{lec01::def:bijective-mapping} \(f:A\to B\):
\[
A\equinum B
\iff
\exists f:A\to B\text{ such that }f\text{ is bijective}.
\]
\end{dfn}

\begin{thm}
  If $\card{A}, \card{B} \in \NN$, ie, $A$ and $B$ are finite sets, then $\card{A} = \card{B} \iff A \equinum B$ 
\end{thm}

\begin{dfn}[Equivalence Relation]{equivalence-relation}
  A relation $\sim$ on a set $A$ is called an \emph{equivalence relation} if it satisfies the following properties: 
  \begin{enumerate}[nosep]
    \item Reflexive: $\forall a \in A, a \sim a$
    \item Symmetric: $\forall a,b \in A, a \sim b \implies b \sim a$
    \item Transitive: $\forall a,b,c \in A, a \sim b, b \sim c \implies a \sim c$
  \end{enumerate}
\end{dfn}

\begin{thm}[Cardinality is an Equivalence Relation]{equivalence-relation-cardinality}
  The Cardinality relation $\equinum$ on the set of all sets is an equivalence relation. 
\end{thm}

\begin{example}
$\NN \equinum \NNOne$. Witness: $f: \NN \to \NNOne, f(n) = n + 1$
\end{example} 

\begin{rem}
For infinite sets, the cardinality maybe still the same even if we remove some (or a finite number of) elements from the set. 
\end{rem} 

\begin{thm}[Integers and Natural Numbers are Equinumerous]{integers-natural-equinum}
  $\ZZ \equinum \NN$. Witness: 
  \begin{align*}
    f = \left\{\begin{array}{ll}
      n/2 & \text{if } n \text{ is even} \\
      -(n+1)/2 & \text{if } n \text{ is odd}
    \end{array}\right.
  \end{align*}
\end{thm}

\begin{thm}[Countability of the Rationals]{rationals-natural-equinum}
  $\QQ \equinum \NN$. Arrange the positive rationals in the Calkin--Wilf tree,
  where each $a/b$ has children $a/(a+b)$ and $(a+b)/b$, and read it
  level by level:
  \[
    \frac11;\qquad
    \frac12,\frac21;\qquad
    \frac13,\frac32,\frac23,\frac31;\qquad \ldots
  \]
  Thus the levels contain $1,2,4,\ldots$ elements. Let
  $r_1,r_2,\ldots$ be this enumeration. Since every positive rational occurs
  exactly once, a witness $f:\NN\to\QQ$ is
  \[
    f(n)=
    \begin{cases}
      0, & n=0,\\
      r_{(n+1)/2}, & n\text{ is odd},\\
      -r_{n/2}, & n\text{ is even and }n>0.
    \end{cases}
  \]
\end{thm}

\begin{thm}[Equinumerosity of Intervals]{interval-equinum}
Every \xref[interval]{def:interval} which is not a single point is equinumerous to another \xref[interval]{def:interval} which is not a single point. 
\end{thm} 

\begin{proof}
Witness: \[f: (a,b) \to (c,d), f(x) = \frac{d-c}{b-a}(x-a)+c\] 
\end{proof} 

\begin{thm}[Intervals and the Real Line]{interval-real-line-equinum}
Every \xref[interval]{def:interval} which is not a single point is equinumerous to $\RR$. 
\end{thm}

\begin{proof}
  For $(-\pi/2, \pi/2)$, the function $f(x) = \tan(x)$ is a bijection from $(-\pi/2, \pi/2)$ to $\RR$. 

  Notice that for every interval $(a,b)$ which is not a single point, we can find a bijection from $(a,b)$ to $(-\pi/2, \pi/2)$ by
  \xref[Theorem~\ref*{thm:interval-equinum}]{thm:interval-equinum}. Since \xref[bijections are transitive]{def:equivalence-relation-cardinality}, we can conclude that every interval $(a,b)$ which is not a single point is equinumerous to $\RR$. 
\end{proof} 

\section{Countability}

\begin{dfn}[Countable Infinite Set]{countable-infinite-set} 
A set $A$ is called \emph{countably infinite} if $N \equinum A$.  
\end{dfn} 

\begin{dfn}
[Countable Set]{countable-set}
A set $A$ is called \emph{countable} if it is either finite or \xref[countably infinite]{def:countable-infinite-set}.  
\end{dfn} 

\begin{thm}[$\ZZ, \QQ$ are Countable]{integers-rationals-countable}
  $\ZZ$ and $\QQ$ are countable since they are equinumerous to $\NN$ according to \xref[Theorem~\ref*{thm:integers-natural-equinum}]{thm:integers-natural-equinum} and \xref[Theorem~\ref*{thm:rationals-natural-equinum}]{thm:rationals-natural-equinum}. 
\end{thm} 

\begin{thm}[Cantor's Theorem]{cantor-theorem}
The interval $(0,1)$ is uncountable. 
\end{thm} 

\begin{proof}
Since $\NN\equinum\NNOne$, suppose there exists a bijection
$f:\NNOne\to(0,1)$. By
\xref[Theorem~\ref*{thm:canonical-decimal-uniqueness}]{thm:canonical-decimal-uniqueness},
choose the unique decimal expansion of each $f(n)$ which is not eventually
equal to $9$. This defines a function
$g:\NNOne\to\NNOne\to\{0,1,2,3,4,5,6,7,8,9\}$.

Then, we can define another function of decimal expansion 
\begin{align}
h(n) = \Upsilon(g(n)(n)),\qquad n\in\NNOne \\ 
\Upsilon = \{(0, 1), (1, 2), (2, 3), (3, 4), (4, 5), (5, 6), (6, 7), (7, 8), (8, 1), (9, 1)\} 
\end{align} 

Then, we can construct
\[
x=\sum_{n=1}^{\infty}h(n)10^{-n}.
\]
Since $h(n)\in\{1,2,\ldots,8\}$, we have $x\in(0,1)$, and this is its
unique canonical decimal expansion by
\xref[Theorem~\ref*{thm:canonical-decimal-uniqueness}]{thm:canonical-decimal-uniqueness}.
For every $n\in\NNOne$, $h(n)\neq g(n)(n)$, so $x\neq f(n)$. Hence $x$ is
not in the image of $f$, contradicting the assumption that $f$ is a
bijection. Therefore, $(0,1)$ is uncountable.
\end{proof}

\begin{thm}[$\RR$ is Uncountable]{reals-uncountable}
  $\RR$ is uncountable. 
\end{thm} 

\begin{proof}
Assume that $\RR$ is countable. According to \xref[Theorem~\ref*{thm:interval-real-line-equinum}]{thm:interval-real-line-equinum}, we can conclude that $(0,1) \equinum \RR \equinum \NN$. According to \xref[Cantor's Theorem]{thm:cantor-theorem}, $(0,1)$ is uncountable, which contradicts the assumption that $\RR$ is countable. Therefore, $\RR$ is uncountable. 
\end{proof} 

\begin{thm}[Countable Union of Countable Sets]{countable-union-of-countable-sets}
  Given $n$ countable sets $A_1, A_2, \ldots, A_n$, the union $\bigcup_{i=1}^{n} A_i$ is countable.
\end{thm} 

\begin{proof}
  We can prove this theorem by induction. 

  Let $A$ and $B$ be two countable sets. If $A$ and $B$ are finite, trivially it is countable. 
\end{proof} 

\section{Power Sets} 

\end{document}
